\documentclass[11pt,a4paper]{article}

\usepackage{kotex}
\usepackage{fontspec}
\setmainfont{Times New Roman}
\setmainhangulfont{AppleMyungjo}
\setsanshangulfont{Apple SD Gothic Neo}

\usepackage[a4paper,top=18mm,bottom=18mm,left=20mm,right=20mm]{geometry}
\usepackage{array}
\usepackage{enumitem}
\linespread{1.18}
\setlength{\parindent}{0pt}

\newcommand{\sol}[2]{\vspace{6pt}{\sffamily\bfseries #1} \quad 정답 \fbox{#2}\par\vspace{1pt}}

\begin{document}

\begin{center}
{\fontsize{15}{17}\selectfont\sffamily\bfseries 2026 고1 영어 변형 — 지문 31 정답 및 해설}\\[2pt]
{\sffamily 변화와 Heraclitus}
\end{center}
\vspace{6pt}

\begin{center}
\renewcommand{\arraystretch}{1.4}
\begin{tabular}{|c|c|c|c|c|}
\hline
\sffamily 문항 & 1 & 2 & 3 & 4 \\\hline
\sffamily 유형 & \sffamily 빈칸추론 & \sffamily 영어 주제 & \sffamily 서술형(해석) & \sffamily 함의추론 \\\hline
\sffamily 정답 & ③ & ④ & (아래) & ③ \\\hline
\sffamily 배점 & 2점 & 2점 & 2점 & 2점 \\\hline
\end{tabular}
\end{center}
\vspace{8pt}\hrule\vspace{4pt}

\sol{1. 빈칸추론}{③}
빈칸은 ``the neural networks in your brain are also \underline{\hspace{2cm}} changing as you perceive each word''에 해당한다.
원문 어휘는 \textit{imperceptibly}(감지할 수 없게)이며, 이는 변화가 일어나고 있지만 인식되지 않는다는 문맥을 나타낸다.
③ \textbf{subtly}(미묘하게, 감지하기 어렵게)가 이 의미와 가장 가깝다.
\par\vspace{2pt}
\textbf{오답 소거:}
① \textit{rapidly}(빠르게) — 본문은 ``a minute amount''로 변화가 미미함을 강조. 속도가 빠르다는 의미와 충돌.
② \textit{intentionally}(의도적으로) — 뇌 신경망의 변화는 의도와 무관한 자동적 과정.
④ \textit{visibly}(눈에 띄게) — 감지할 수 없다는 문맥과 반대.
⑤ \textit{permanently}(영구적으로) — 변화의 지속성이 아니라 비가시성이 핵심.

\sol{2. 영어 주제}{④}
글 전체는 우리 세계가 끊임없이 변화하며, 그 변화 대부분이 감지되지 않는다는 주장을 전개한다.
Heraclitus의 인용과 Cratylus의 보완 모두 ``변화는 항상 일어난다''는 핵심 메시지를 뒷받침한다.
④ \textbf{the continuous and often unnoticed nature of change in our world}가 이를 정확히 요약한다.
\par\vspace{2pt}
\textbf{오답 소거:}
① 두 철학자의 \textit{논쟁(debate)}은 없다. Cratylus는 Heraclitus의 견해를 \textit{확장}한 것.
② 신경망의 언어 처리는 지엽적 예시일 뿐, 글 전체 주제가 아님.
③ 마음 챙김(mindfulness) 개념은 언급되지 않음.
⑤ 현대 과학과 고대 철학의 확인 관계는 글의 초점이 아님.

\vspace{6pt}{\sffamily\bfseries 3. 서술형 — 해석}\par\vspace{2pt}

\noindent\fbox{대상 문장}\par
\textit{Crucially, even when we're seemingly not doing anything of note, events are taking place outside your immediate surroundings that will change your life in the future, though you won't realize it yet.}

\par\vspace{6pt}
\noindent\fbox{모범 답안}\par
\textbf{결정적으로, 우리가 특별히 아무것도 하지 않는 것처럼 보일 때조차, 미래에 당신의 삶을 바꿀 사건들이 당신의 즉각적인 주변 환경 밖에서 일어나고 있다. 비록 당신이 아직 그것을 깨닫지 못하더라도.}

\par\vspace{4pt}
\noindent\textbf{채점 포인트 (각 1점씩, 총 2점):}
\begin{itemize}[leftmargin=1.5em,itemsep=1pt,topsep=2pt]
  \item \textit{seemingly not doing anything of note} → ``특별히 아무것도 하지 않는 것처럼 보일 때'' (또는 동의어)
  \item \textit{outside your immediate surroundings} → ``즉각적인[바로 주변의] 환경 밖에서''
  \item \textit{though you won't realize it yet} → ``비록 아직 깨닫지 못하더라도'' (또는 동의어)
\end{itemize}
(세 포인트 중 2개 이상 포함 시 2점, 1개 포함 시 1점, 0개 시 0점)

\vspace{8pt}
\sol{4. 함의 추론}{③}
Heraclitus의 강물 비유는 같은 강에 다시 들어갈 수 없다는 말로, 강도 변하고 그 강에 들어가는 사람도 변한다는 뜻이다. 이는 세계와 관찰자 모두가 계속 변화한다는 글의 핵심과 연결되므로 ③이 정답이다.

\end{document}
